%% Second example research chapter.  This one keeps all content in a
%% single file -- fine for shorter chapters.  For longer,
%% paper-length chapters, follow the per-section file layout of
%% chapter-example-1/ instead.

% Chapter-specific commands
\newcommand{\tooltwo}{\textsc{ToolTwo}\xspace}

\chapter{Title of Your Second Research Chapter}
\label{chap:example-two}

% This chapter is based on:
%   Your Name, Coauthor One, and Coauthor Two. "Title of the Second
%   Paper". Venue Name (VENUE 20XX).

\section{Introduction}

Replace this text with the content of your second research chapter.
This chapter builds on \autoref{chap:example-one} and presents
\tooltwo.

\section{Approach}

Describe the approach.

\section{Evaluation}

Describe the evaluation.

\section{Conclusion}

Summarize the chapter.
